\subsection{\label{sec.gf.lim}Limits of FPSs}

Much of what we have been doing with FPSs has been an analogue (a mockery, as
some might say) of analysis, particularly of complex analysis. We have defined
FPSs, derivatives, infinite sums and products, exponentials and various
related concepts, all by mimicking the respective analytic notions but
avoiding any reference to real numbers or any other peculiarities that our
base ring $K$ may or may not have.

Let me introduce yet another such formal analogue of a classical concept from
analysis: that of a limit. Specifically, I will define (coefficientwise)
limits of sequences of FPSs. The exposition will follow \cite[\S 7.5]%
{Loehr-BC}. These limits are very easy to define, yet quite useful. (In
particular, they can be used to give a simpler definition of infinite
products, although a less natural one.)

\subsubsection{Stabilization of scalars}

We start with the most trivial kind of limit whatsoever: the limit of an
eventually constant sequence. Its main advantage is its generality and the
simplicity of its definition:

\begin{definition}
\label{def.fps.lim.stab}Let $\left(  a_{i}\right)  _{i\in\mathbb{N}}=\left(
a_{0},a_{1},a_{2},\ldots\right)  \in K^{\mathbb{N}}$ be a sequence of elements
of $K$. Let $a\in K$.

We say that the sequence $\left(  a_{i}\right)  _{i\in\mathbb{N}}$
\emph{stabilizes to }$a$ if there exists some $N\in\mathbb{N}$ such that%
\[
\text{all integers }i\geq N\text{ satisfy }a_{i}=a.
\]


Instead of saying \textquotedblleft$\left(  a_{i}\right)  _{i\in\mathbb{N}}$
stabilizes to $a$\textquotedblright, we can also say \textquotedblleft$a_{i}$
stabilizes to $a$ as $i\rightarrow\infty$\textquotedblright. (The words
\textquotedblleft as $i\rightarrow\infty$\textquotedblright\ here signify that
we are talking not about a single term $a_{i}$ but about the entire sequence
$\left(  a_{i}\right)  _{i\in\mathbb{N}}$.)

If $a_{i}$ stabilizes to $a$ as $i\rightarrow\infty$, then we write
$\lim\limits_{i\rightarrow\infty}a_{i}=a$ and say that $a$ is the \emph{limit}
(or \emph{eventual value}) of $\left(  a_{i}\right)  _{i\in\mathbb{N}}$ (or,
less precisely, that $a$ is the limit of the $a_{i}$). This is legitimate,
since $a$ is uniquely determined by the sequence $\left(  a_{i}\right)
_{i\in\mathbb{N}}$.

We can replace $\mathbb{N}$ by $\mathbb{Z}_{\geq q}=\left\{  q,q+1,q+2,\ldots
\right\}  $ in this definition, where $q$ is any integer. That is, everything
we said applies not just to sequences of the form $\left(  a_{i}\right)
_{i\in\mathbb{N}}=\left(  a_{0},a_{1},a_{2},\ldots\right)  $, but also to
sequences of the form $\left(  a_{i}\right)  _{i\geq q}=\left(  a_{q}%
,a_{q+1},a_{q+2},\ldots\right)  $.
\end{definition}

Here are some examples and non-examples of stabilization:

\begin{itemize}
\item The sequence%
\[
\left(  \left\lfloor \dfrac{5}{i}\right\rfloor \right)  _{i\geq1}=\left(
\left\lfloor \dfrac{5}{1}\right\rfloor ,\left\lfloor \dfrac{5}{2}\right\rfloor
,\left\lfloor \dfrac{5}{3}\right\rfloor ,\ldots\right)  =\left(
5,2,1,1,1,0,0,0,0,\ldots\right)  \in\mathbb{Z}^{\mathbb{N}}%
\]
stabilizes to $0$, since all integers $i\geq6$ satisfy $\left\lfloor \dfrac
{5}{i}\right\rfloor =0$.

\item On the other hand, the sequence $\left(  \dfrac{5}{i}\right)  _{i\geq
1}\in\mathbb{Z}^{\mathbb{N}}$ does not stabilize to anything. (It does
converge to $0$ in the sense of real analysis, but this is a weaker statement
than stabilization.)

\item Any constant sequence $\left(  a,a,a,\ldots\right)  \in K^{\mathbb{N}}$
stabilizes to $a$.

\item A well-known (if non-constructive) fact says that every weakly
decreasing sequence of nonnegative integers stabilizes to some nonnegative
integer (since it cannot decrease infinitely often).

\item If $s$ is a nilpotent element of a ring $K$ (that is, an element
satisfying $s^{k}=0$ for some $k\geq0$), then the sequence $\left(
s^{i}\right)  _{i\in\mathbb{N}}=\left(  s^{0},s^{1},s^{2},\ldots\right)  $
stabilizes to $0\in K$. For example, the element $\overline{2}$ of the ring
$\mathbb{Z}/8\mathbb{Z}$ is nilpotent (with $\overline{2}^{3}=\overline{0}$),
so that the sequence $\left(  \overline{2}^{i}\right)  _{i\in\mathbb{N}%
}=\left(  \overline{2},\overline{4},\overline{0},\overline{0},\overline
{0},\ldots\right)  $ stabilizes to $\overline{0}$.

\item If $s$ is an idempotent element of a ring $K$ (that is, an element
satisfying $s^{2}=s$), then the sequence $\left(  s^{i}\right)  _{i\in
\mathbb{N}}=\left(  s^{0},s^{1},s^{2},\ldots\right)  $ stabilizes to $s$ (in
fact, $s^{1}=s^{2}=s^{3}=\cdots=s$).
\end{itemize}

\begin{remark}
If you are familiar with point-set topology, then you will recognize
Definition \ref{def.fps.lim.stab} as an instance of topological convergence:
Namely, a sequence $\left(  a_{i}\right)  _{i\in\mathbb{N}}$ stabilizes to
some $a\in K$ if and only if $\left(  a_{i}\right)  _{i\in\mathbb{N}}$
converges to $a$ in the discrete topology.
\end{remark}

\subsubsection{Coefficientwise stabilization of FPSs}

Now we define a somewhat subtler notion of limits, one specific to FPSs (as
opposed to elements of an arbitrary commutative ring). This is the notion that
we will use in what follows:

\begin{definition}
\label{def.fps.lim.coeff-stab}Let $\left(  f_{i}\right)  _{i\in\mathbb{N}}\in
K\left[  \left[  x\right]  \right]  ^{\mathbb{N}}$ be a sequence of FPSs over
$K$. Let $f\in K\left[  \left[  x\right]  \right]  $ be an FPS.

We say that $\left(  f_{i}\right)  _{i\in\mathbb{N}}$ \emph{coefficientwise
stabilizes to }$f$ if for each $n\in\mathbb{N}$,%
\[
\text{the sequence }\left(  \left[  x^{n}\right]  f_{i}\right)  _{i\in
\mathbb{N}}\text{ stabilizes to }\left[  x^{n}\right]  f.
\]


Instead of saying \textquotedblleft$\left(  f_{i}\right)  _{i\in\mathbb{N}}$
coefficientwise stabilizes to $f$\textquotedblright, we can also say
\textquotedblleft$f_{i}$ coefficientwise stabilizes to $f$ as $i\rightarrow
\infty$\textquotedblright\ or write \textquotedblleft$f_{i}\rightarrow f$ as
$i\rightarrow\infty$\textquotedblright.

If $f_{i}$ coefficientwise stabilizes to $f$ as $i\rightarrow\infty$, then we
write $\lim\limits_{i\rightarrow\infty}f_{i}=f$ and say that $f$ is the
\emph{limit} of $\left(  f_{i}\right)  _{i\in\mathbb{N}}$ (or, less precisely,
that $f$ is the limit of the $f_{i}$). This is legitimate, because $f$ is
uniquely determined by the sequence $\left(  f_{i}\right)  _{i\in\mathbb{N}}$.

Again, we can replace $\mathbb{N}$ by $\mathbb{Z}_{\geq q}=\left\{
q,q+1,q+2,\ldots\right\}  $ in this definition, where $q$ is any integer.
\end{definition}

\begin{example}
\textbf{(a)} The sequence $\left(  x^{i}\right)  _{i\in\mathbb{N}}$ of FPSs
coefficientwise stabilizes to the FPS $0$. (Indeed, for each $n\in\mathbb{N}$,
the sequence $\left(  \left[  x^{n}\right]  \left(  x^{i}\right)  \right)
_{i\in\mathbb{N}}$ consists of a single $1$ and infinitely many $0$s; thus,
this sequence stabilizes to $0$.) In other words, we have $\lim
\limits_{i\rightarrow\infty}x^{i}=0$. \medskip

\textbf{(b)} The sequence $\left(  \dfrac{1}{i}x\right)  _{i\geq1}$ of FPSs
does not stabilize (since the $x^{1}$-coefficients of these FPSs $\dfrac{1}%
{i}x$ never stabilize). Thus, $\lim\limits_{i\rightarrow\infty}\left(
\dfrac{1}{i}x\right)  $ does not exist. \medskip

\textbf{(c)} We have%
\[
\lim\limits_{i\rightarrow\infty}\left(  \left(  1+x^{1}\right)  \left(
1+x^{2}\right)  \cdots\left(  1+x^{i}\right)  \right)  =\prod_{k=1}^{\infty
}\left(  1+x^{k}\right)  .
\]
Indeed, the $x^{n}$-coefficient of the product $\left(  1+x^{1}\right)
\left(  1+x^{2}\right)  \cdots\left(  1+x^{i}\right)  $ depends only on those
factors in which the exponent is $\leq n$, and thus stops changing after $i$
surpasses $n$; therefore, it stabilizes to the $x^{n}$-coefficient of the
infinite product $\prod_{k=1}^{\infty}\left(  1+x^{k}\right)  $. \medskip

\textbf{(d)} It would be nice to have $\lim\limits_{n\rightarrow\infty}\left(
1+\dfrac{x}{n}\right)  ^{n}=\exp$, as in real analysis. Unfortunately, this is
not the case. In fact, the binomial formula yields
\[
\left(  1+\dfrac{x}{n}\right)  ^{n}=1+\underbrace{n\cdot\dfrac{x}{n}}%
_{=x}+\,\dbinom{n}{2}\cdot\left(  \dfrac{x}{n}\right)  ^{2}+\cdots
=1+x+\dfrac{\dbinom{n}{2}}{n^{2}}x^{2}+\cdots.
\]
This shows that the $x^{0}$-coefficient and the $x^{1}$-coefficient stabilize
of $\left(  1+\dfrac{x}{n}\right)  ^{n}$ as $n\rightarrow\infty$, but the
$x^{2}$-coefficient does not. Thus, $\lim\limits_{n\rightarrow\infty}\left(
1+\dfrac{x}{n}\right)  ^{n}$ does not exist (according to our definition of
limit). \medskip

\textbf{(e)} Coefficientwise stabilization is weaker than stabilization: If a
sequence $\left(  f_{0},f_{1},f_{2},\ldots\right)  $ of FPSs stabilizes to an
FPS $f$ (meaning that $f_{i}=f$ for all sufficiently high $i\in\mathbb{N}$),
then it coefficientwise stabilizes to $f$ as well. In particular, any constant
sequence $\left(  f,f,f,\ldots\right)  $ of FPSs stabilizes to $f$.
\end{example}

\begin{remark}
If you are familiar with point-set topology, then you will again recognize
Definition \ref{def.fps.lim.coeff-stab} as an instance of topological
convergence: Namely, we recall that each FPS is an infinite sequence of
elements of $K$ (its coefficients). Thus, the set $K\left[  \left[  x\right]
\right]  $ is the infinite Cartesian product $K\times K\times K\times\cdots$.
If we equip each factor $K$ in this product with the discrete topology, then
the entire product $K\left[  \left[  x\right]  \right]  $ becomes a product
space equipped with the product topology. Now, a sequence $\left(
f_{i}\right)  _{i\in\mathbb{N}}$ of FPSs in $K\left[  \left[  x\right]
\right]  $ coefficientwise stabilizes to some FPS $f\in K\left[  \left[
x\right]  \right]  $ if and only if $\left(  f_{i}\right)  _{i\in\mathbb{N}}$
converges to $f$ in this product topology. (This is not surprising: After all,
the product topology is also known as the \textquotedblleft topology of
pointwise convergence\textquotedblright, and in our case \textquotedblleft
pointwise\textquotedblright\ means \textquotedblleft
coefficientwise\textquotedblright.)
\end{remark}

\subsubsection{Some properties of limits}

We shall now state some basic properties of limits of FPSs. All the proofs are
simple and therefore omitted; some can be found in the Appendix (Section
\ref{sec.details.gf.lim}).

\begin{theorem}
\label{thm.fps.lim.lim-crit}Let $\left(  f_{i}\right)  _{i\in\mathbb{N}}\in
K\left[  \left[  x\right]  \right]  ^{\mathbb{N}}$ be a sequence of FPSs.
Assume that for each $n\in\mathbb{N}$, there exists some $g_{n}\in K$ such
that the sequence $\left(  \left[  x^{n}\right]  f_{i}\right)  _{i\in
\mathbb{N}}$ stabilizes to $g_{n}$. Then, the sequence $\left(  f_{i}\right)
_{i\in\mathbb{N}}$ coefficientwise stabilizes to $\sum_{n\in\mathbb{N}}%
g_{n}x^{n}$. (That is, $\lim\limits_{i\rightarrow\infty}f_{i}=\sum
_{n\in\mathbb{N}}g_{n}x^{n}$.)
\end{theorem}

\begin{proof}
Obvious.
\end{proof}

The following easy lemma connects limits with the notion of $x^{n}%
$-equivalence (as introduced in Definition \ref{def.fps.xneq}):

\begin{lemma}
\label{lem.fps.lim.xn-equiv}Assume that $\left(  f_{i}\right)  _{i\in
\mathbb{N}}$ is a sequence of FPSs, and that $f$ is an FPS such that
$\lim\limits_{i\rightarrow\infty}f_{i}=f$. Then, for each $n\in\mathbb{N}$,
there exists some integer $N\in\mathbb{N}$ such that
\[
\text{all integers }i\geq N\text{ satisfy }f_{i}\overset{x^{n}}{\equiv}f.
\]

\end{lemma}

\begin{proof}
[Proof of Lemma \ref{lem.fps.lim.xn-equiv} (sketched).]Let $n\in\mathbb{N}$.
For each $k\in\left\{  0,1,\ldots,n\right\}  $, we pick an $N_{k}\in
\mathbb{N}$ such that all $i\geq N_{k}$ satisfy $\left[  x^{k}\right]
f_{i}=\left[  x^{k}\right]  f$ (such an $N_{k}$ exists, since $\lim
\limits_{i\rightarrow\infty}f_{i}=f$). Then, all integers $i\geq\max\left\{
N_{0},N_{1},\ldots,N_{n}\right\}  $ satisfy $f_{i}\overset{x^{n}}{\equiv}f$.
See Section \ref{sec.details.gf.lim} for the details of this proof.
\end{proof}

The following proposition is an analogue of the classical \textquotedblleft
limits respect sums and products\textquotedblright\ theorem from real analysis:

\begin{proposition}
\label{prop.fps.lim.sum-prod}Assume that $\left(  f_{i}\right)  _{i\in
\mathbb{N}}$ and $\left(  g_{i}\right)  _{i\in\mathbb{N}}$ are two sequences
of FPSs, and that $f$ and $g$ are two FPSs such that%
\[
\lim\limits_{i\rightarrow\infty}f_{i}=f\ \ \ \ \ \ \ \ \ \ \text{and}%
\ \ \ \ \ \ \ \ \ \ \lim\limits_{i\rightarrow\infty}g_{i}=g.
\]


Then,
\[
\lim\limits_{i\rightarrow\infty}\left(  f_{i}+g_{i}\right)
=f+g\ \ \ \ \ \ \ \ \ \ \text{and}\ \ \ \ \ \ \ \ \ \ \lim
\limits_{i\rightarrow\infty}\left(  f_{i}g_{i}\right)  =fg.
\]
(In other words, the sequences $\left(  f_{i}+g_{i}\right)  _{i\in\mathbb{N}}$
and $\left(  f_{i}g_{i}\right)  _{i\in\mathbb{N}}$ coefficientwise stabilize
to $f+g$ and $fg$, respectively.)
\end{proposition}

\begin{proof}
[Proof of Proposition \ref{prop.fps.lim.sum-prod} (sketched).]Let
$n\in\mathbb{N}$. Then, Lemma \ref{lem.fps.lim.xn-equiv} shows that there
exists some integer $N\in\mathbb{N}$ such that
\[
\text{all integers }i\geq N\text{ satisfy }f_{i}\overset{x^{n}}{\equiv}f.
\]
Similarly, there exists some integer $M\in\mathbb{N}$ such that%
\[
\text{all integers }i\geq M\text{ satisfy }g_{i}\overset{x^{n}}{\equiv}g.
\]
Pick such integers $N$ and $M$, and set $P:=\max\left\{  N,M\right\}  $. Then,
all integers $i\geq P$ satisfy both $f_{i}\overset{x^{n}}{\equiv}f$ and
$g_{i}\overset{x^{n}}{\equiv}g$, and therefore $f_{i}g_{i}\overset{x^{n}%
}{\equiv}fg$ (by (\ref{eq.thm.fps.xneq.props.b.*})), and thus $\left[
x^{n}\right]  \left(  f_{i}g_{i}\right)  =\left[  x^{n}\right]  \left(
fg\right)  $. This shows that the sequence $\left(  \left[  x^{n}\right]
\left(  f_{i}g_{i}\right)  \right)  _{i\in\mathbb{N}}$ stabilizes to $\left[
x^{n}\right]  \left(  fg\right)  $. Since this holds for each $n\in\mathbb{N}%
$, we conclude that $\lim\limits_{i\rightarrow\infty}\left(  f_{i}%
g_{i}\right)  =fg$. The proof of $\lim\limits_{i\rightarrow\infty}\left(
f_{i}+g_{i}\right)  =f+g$ is analogous.

See Section \ref{sec.details.gf.lim} for the details of this proof.
\end{proof}

\begin{corollary}
\label{cor.fps.lim.sum-prod-k}Let $k\in\mathbb{N}$. For each $i\in\left\{
1,2,\ldots,k\right\}  $, let $f_{i}$ be an FPS, and let $\left(
f_{i,n}\right)  _{n\in\mathbb{N}}$ be a sequence of FPSs such that%
\[
\lim\limits_{n\rightarrow\infty}\left(  f_{i,n}\right)  =f_{i}%
\]
(note that it is $n$, not $i$, that goes to $\infty$ here!). Then,%
\[
\lim\limits_{n\rightarrow\infty}\sum_{i=1}^{k}f_{i,n}=\sum_{i=1}^{k}%
f_{i}\ \ \ \ \ \ \ \ \ \ \text{and}\ \ \ \ \ \ \ \ \ \ \lim
\limits_{n\rightarrow\infty}\prod_{i=1}^{k}f_{i,n}=\prod_{i=1}^{k}f_{i}.
\]

\end{corollary}

\begin{proof}
Follows by induction on $k$, using Proposition \ref{prop.fps.lim.sum-prod}.
\end{proof}

\begin{remark}
The analogue of Corollary \ref{cor.fps.lim.sum-prod-k} for infinite sums (or
products) is not true without further requirements. For instance, for each
$i\in\mathbb{N}$, we can define a constant FPS $f_{i,n}=\delta_{i,n}$ (using
Definition \ref{def.kron-delta}). Then, we have $\lim\limits_{n\rightarrow
\infty}f_{i,n}=0$ for any fixed $i\in\mathbb{N}$, but $\lim
\limits_{n\rightarrow\infty}\underbrace{\sum_{i=0}^{\infty}f_{i,n}}_{=1}%
=\lim\limits_{n\rightarrow\infty}1=1$.
\end{remark}

The following proposition says that limits respect quotients:

\begin{proposition}
\label{prop.fps.lim.sum-quot}Assume that $\left(  f_{i}\right)  _{i\in
\mathbb{N}}$ and $\left(  g_{i}\right)  _{i\in\mathbb{N}}$ are two sequences
of FPSs, and that $f$ and $g$ are two FPSs such that%
\[
\lim\limits_{i\rightarrow\infty}f_{i}=f\ \ \ \ \ \ \ \ \ \ \text{and}%
\ \ \ \ \ \ \ \ \ \ \lim\limits_{i\rightarrow\infty}g_{i}=g.
\]
Assume that each FPS $g_{i}$ is invertible. Then, $g$ is also invertible, and
we have%
\[
\lim\limits_{i\rightarrow\infty}\dfrac{f_{i}}{g_{i}}=\dfrac{f}{g}.
\]

\end{proposition}

\begin{proof}
[Proof of Proposition \ref{prop.fps.lim.sum-quot} (sketched).]First use
Proposition \ref{prop.fps.invertible} to show that $g$ is invertible; then use
Theorem \ref{thm.fps.xneq.props} \textbf{(e)} to prove $\lim
\limits_{i\rightarrow\infty}\dfrac{f_{i}}{g_{i}}=\dfrac{f}{g}$. A detailed
proof can be found in Section \ref{sec.details.gf.lim}.
\end{proof}

Limits of FPSs furthermore respect composition:

\begin{proposition}
\label{prop.fps.lim.comp}Assume that $\left(  f_{i}\right)  _{i\in\mathbb{N}}$
and $\left(  g_{i}\right)  _{i\in\mathbb{N}}$ are two sequences of FPSs, and
that $f$ and $g$ are two FPSs such that%
\[
\lim\limits_{i\rightarrow\infty}f_{i}=f\ \ \ \ \ \ \ \ \ \ \text{and}%
\ \ \ \ \ \ \ \ \ \ \lim\limits_{i\rightarrow\infty}g_{i}=g.
\]
Assume that $\left[  x^{0}\right]  g_{i}=0$ for each $i\in\mathbb{N}$. Then,
$\left[  x^{0}\right]  g=0$ and%
\[
\lim\limits_{i\rightarrow\infty}\left(  f_{i}\circ g_{i}\right)  =f\circ g.
\]

\end{proposition}

\begin{proof}
[Proof of Proposition \ref{prop.fps.lim.comp} (sketched).]Similar to
Proposition \ref{prop.fps.lim.sum-quot}, but now using Proposition
\ref{prop.fps.xneq.comp}. See Exercise \ref{exe.fps.lim.sum-prod-quot}
\textbf{(b)} for details.
\end{proof}

Limits of FPSs also respect derivatives (unlike in real analysis, where this
\href{https://en.wikipedia.org/wiki/Uniform_convergence#To_differentiability}{holds
only under subtle additional conditions}):

\begin{proposition}
\label{prop.fps.lim.deriv-lim}Let $\left(  f_{n}\right)  _{n\in\mathbb{N}}$ be
a sequence of FPSs, and let $f$ be an FPS such that%
\[
\lim\limits_{i\rightarrow\infty}f_{i}=f.
\]
Then,%
\[
\lim\limits_{i\rightarrow\infty}f_{i}^{\prime}=f^{\prime}.
\]

\end{proposition}

\begin{proof}
This is nearly trivial (see Exercise \ref{exe.fps.lim.sum-prod-quot}
\textbf{(a)}).
\end{proof}

Next, let us see how infinite sums and infinite products can be written as
limits of finite sums and products, at least as long as they are indexed by
$\mathbb{N}$. The following two theorems and one corollary are easy to prove;
detailed proofs can be found in Section \ref{sec.details.gf.lim}.

\begin{theorem}
\label{thm.fps.lim.sum-lim}Let $\left(  f_{n}\right)  _{n\in\mathbb{N}}$ be a
summable sequence of FPSs. Then,%
\[
\lim\limits_{i\rightarrow\infty}\sum_{n=0}^{i}f_{n}=\sum_{n\in\mathbb{N}}%
f_{n}.
\]
In other words, the infinite sum $\sum_{n\in\mathbb{N}}f_{n}$ is the limit of
the finite partial sums $\sum_{n=0}^{i}f_{n}$.
\end{theorem}

\begin{theorem}
\label{thm.fps.lim.prod-lim}Let $\left(  f_{n}\right)  _{n\in\mathbb{N}}$ be a
multipliable sequence of FPSs. Then,%
\[
\lim\limits_{i\rightarrow\infty}\prod_{n=0}^{i}f_{n}=\prod_{n\in\mathbb{N}%
}f_{n}.
\]
In other words, the infinite product $\prod_{n\in\mathbb{N}}f_{n}$ is the
limit of the finite partial products $\prod_{n=0}^{i}f_{n}$.
\end{theorem}

\begin{corollary}
\label{cor.fps.lim.fps-as-pol}Each FPS is a limit of a sequence of
polynomials. Indeed, if $a=\sum_{n\in\mathbb{N}}a_{n}x^{n}$ (with $a_{n}\in
K$), then%
\[
a=\lim\limits_{i\rightarrow\infty}\sum_{n=0}^{i}a_{n}x^{n}.
\]

\end{corollary}

This corollary can be restated as \textquotedblleft the polynomials are dense
in the FPSs\textquotedblright\ (more formally: $K\left[  x\right]  $ is dense
in $K\left[  \left[  x\right]  \right]  $). This fact is useful, as it allows
you to restrict yourself to polynomials when proving some properties of FPSs.
For example, if you want to prove that some identity (such as the Leibniz rule
$\left(  fg\right)  ^{\prime}=f^{\prime}g+fg^{\prime}$) holds for all FPSs, it
suffices to prove it for polynomials, and then conclude the general case using
a limiting argument.

Both Theorem \ref{thm.fps.lim.sum-lim} and Theorem \ref{thm.fps.lim.prod-lim}
have converses (which are somewhat harder to prove):

\begin{theorem}
\label{thm.fps.lim.sum-lim-conv}Let $\left(  f_{0},f_{1},f_{2},\ldots\right)
$ be a sequence of FPSs such that $\lim\limits_{i\rightarrow\infty}\sum
_{n=0}^{i}f_{n}$ exists. Then, the family $\left(  f_{n}\right)
_{n\in\mathbb{N}}$ is summable, and satisfies%
\[
\sum_{n\in\mathbb{N}}f_{n}=\lim\limits_{i\rightarrow\infty}\sum_{n=0}^{i}%
f_{n}.
\]

\end{theorem}

\begin{theorem}
\label{thm.fps.lim.prod-lim-conv}Let $\left(  f_{0},f_{1},f_{2},\ldots\right)
$ be a sequence of FPSs such that $\lim\limits_{i\rightarrow\infty}\prod
_{n=0}^{i}f_{n}$ exists. Then, the family $\left(  f_{n}\right)
_{n\in\mathbb{N}}$ is multipliable, and satisfies%
\[
\prod_{n\in\mathbb{N}}f_{n}=\lim\limits_{i\rightarrow\infty}\prod_{n=0}%
^{i}f_{n}.
\]

\end{theorem}

We refer to Section \ref{sec.details.gf.lim} for the proofs of these two theorems.

Theorem \ref{thm.fps.lim.sum-lim} and Theorem \ref{thm.fps.lim.sum-lim-conv}
combined show that we could have just as well defined infinite sums of FPSs as
limits of finite partial sums, as long as we were happy with sums of the form
$\sum_{n\in\mathbb{N}}f_{n}$ (as opposed to sums of the form $\sum_{\left(
n,m\right)  \in\mathbb{N}\times\mathbb{N}}$ or other kinds). Likewise, Theorem
\ref{thm.fps.lim.prod-lim} and Theorem \ref{thm.fps.lim.prod-lim-conv} show
the same for products. However, I consider my original definitions to be more
natural. \medskip

It should be clear that all the above-stated properties of limits remain true
if the sequences involved are no longer indexed by $\mathbb{N}$ but rather
indexed by $\mathbb{Z}_{\geq q}=\left\{  q,q+1,q+2,\ldots\right\}  $ for some
integer $q$. Even better: The limit of a sequence $\left(  f_{n}\right)
_{n\geq q}$ does not depend on any chosen finite piece of this sequence, so it
is simultaneously the limit of the sequence $\left(  f_{n}\right)  _{n\geq r}$
for any integer $r$ (as long as $f_{n}$ is well-defined for all $n\geq r$). In
this regard, limits behave exactly as in real analysis.
